
\chapter{章末問題の略解}

\begin{answerof}{prob:1:1}
  サイズ(型)は$(2,3)$.
  $A$の$(1,2)$-成分は$2$.
  $A$の$2$-行目は
  \begin{align*}\begin{pmatrix}1&2&3\end{pmatrix}.\end{align*}
  $A$の$2$-列目は
  \begin{align*}\begin{pmatrix}2\\7\end{pmatrix}.\end{align*}
  $A$の転置は
  \begin{align*}\begin{pmatrix}1&6\\2&7\\3&8\end{pmatrix}.\end{align*}
\end{answerof}

\begin{answerof}{prob:1:2}
  具体的に成分を書き下せばわかる.
  \begin{enumerate}
  \item
    $(i,j)=(1,1)$のとき, $5i+j-5=1$.
    $(i,j)=(1,2)$のとき, $5i+j-5=2$.
    $(i,j)=(2,1)$のとき, $5i+j-5=6$.
    $(i,j)=(2,2)$のとき, $5i+j-5=7$.
    したがって,
    \begin{align*}A=\begin{pmatrix}1&2\\6&7\end{pmatrix}.\end{align*}
  \item
    $i\neq j$のとき, $2^i\delta_{i,j}=2^i\cdot 0=0$.
    $i=j=1$のとき, $2^i\delta_{i,j}=2^1\cdot 1=2$.
    $i=j=2$のとき, $2^i\delta_{i,j}=2^2\cdot 1=4$.
    したがって,
    \begin{align*}A=\begin{pmatrix}2&0\\0&4\end{pmatrix}.\end{align*}
  \item
    $i\neq 2$のとき, $\delta_{i,2}\delta_{j,1}=0\cdot \delta_{j,1}=0$.
    $j\neq 1$のとき, $\delta_{i,2}\delta_{j,1}=\delta_{i,2}\cdot 0=0$.
    $i=2, j=1$のとき, $\delta_{i,2}\delta_{j,1}=1\cdot 1=1$.
    したがって,
    \begin{align*}A=\begin{pmatrix}0&0\\1&0\end{pmatrix}.\end{align*}
  \end{enumerate}
\end{answerof}
\begin{answerof}{prob:1:3}
  \begin{align*}A&=
    \begin{pmatrix}
      a+1& 3\\
      4+c &5d-10
    \end{pmatrix}\\
    B&=
    \begin{pmatrix}
      7& 2b+5\\
      6 &-4d+8
    \end{pmatrix}
  \end{align*}
  とする.
  行列が等しいということは, 対応する各成分が等しいということなので,
  $A=B$とすると
  \begin{align*}
    \begin{cases}
      a+1=7\\
      3=2b+5\\
      4+c=6\\
      5d-10=-4d+8
    \end{cases}
  \end{align*}
  となる.
  したがって,
  \begin{align*}
    \begin{cases}
      a=6\\
      b=-1\\
      c=2\\
      d=2
    \end{cases}
  \end{align*}
  を得る.
  実際, $(a,b,c,d)=(6,-1,2,2)$とすると,
  \begin{align*}A&=
    \begin{pmatrix}
      a+1& 3\\
      4+c &5d-10
    \end{pmatrix}
    =
    \begin{pmatrix}
      6+1& 3\\
      4+2 &5\cdot 2-10
    \end{pmatrix}
    =
    \begin{pmatrix}
      7& 3\\
      6&0
    \end{pmatrix}
    \\
    B&=
    \begin{pmatrix}
      7& 2b+5\\
      6 &-4d+8
    \end{pmatrix}
=    \begin{pmatrix}
      7& 2(-1)+5\\
      6 &-4\cdot 2+8
    \end{pmatrix}
=    \begin{pmatrix}
      7& 3\\
      6 &0
    \end{pmatrix}
  \end{align*}
  である.
\end{answerof}

\begin{answerof}{prob:1:4}
  \begin{align*}
    A=
    \begin{pmatrix}
      1& 3\\
      4+a &5
    \end{pmatrix}
  \end{align*}
  とする.  $A$が対称行列であるということは, $A=\transposed{A}$ということである.
  \begin{align*}
    \transposed{A}=
    \begin{pmatrix}
      1& 4+a\\
      3 &5
    \end{pmatrix}
  \end{align*}
  であるので, $A=\transposed{A}$とすると,
  \begin{align*}
    \begin{pmatrix}
      1& 3\\
      4+a &5
    \end{pmatrix}
=
    \begin{pmatrix}
      1& 4+a\\
      3 &5
    \end{pmatrix}
  \end{align*}
  である.
  したがって, 
  \begin{align*}
    \begin{cases}
      1=1\\
      3=4+a\\
      4+a=3\\
      5=5
    \end{cases}
  \end{align*}
  となり, $a=-1$を得る.
  実際$a=-1$のとき,
  \begin{align*}
    A&=
    \begin{pmatrix}
      1& 3\\
      4+a &5
    \end{pmatrix}=
    \begin{pmatrix}
      1& 3\\
      3 &5
    \end{pmatrix}\\
    \transposed{A}&=
    \begin{pmatrix}
      1& 3\\
      3 &5
    \end{pmatrix}
  \end{align*}
  であるので対称行列である.
\end{answerof}


\begin{answerof}{prob:1:5}
  \begin{align*}
    A=
    \begin{pmatrix}
      1-a& 3+b\\
      4  &0
    \end{pmatrix}
  \end{align*}
  とする.
  $A$が交代行列であるということは, $-A=\transposed{A}$ということである.
  \begin{align*}
    -A&=
    -\begin{pmatrix}
      1-a& 3+b\\
      4  &0
    \end{pmatrix}
=    \begin{pmatrix}
      -1+a& -3-b\\
      -4  &0
    \end{pmatrix}\\    
    \transposed{A}&=
    \begin{pmatrix}
      1-a& 4\\
      3+b &0
    \end{pmatrix}
  \end{align*}
  であるので, $-A=\transposed{A}$とすると,
  \begin{align*}
    \begin{pmatrix}
      -1+a& -3-b\\
      -4  &0
    \end{pmatrix}&=
    \begin{pmatrix}
      1-a& 4\\
      3+b &0
    \end{pmatrix}
  \end{align*}
  である.
  したがって, 
  \begin{align*}
    \begin{cases}
      -1+a=1-a\\
      -3-b=4\\
      -4=3+b\\
      0=0
    \end{cases}
  \end{align*}
  となり, $a=1$, $b=-7$を得る.
  実際$(a,b)=(1,-7)$のとき,
  \begin{align*}
    -A&=
    \begin{pmatrix}
      -1+a& -3-b\\
      -4  &0
    \end{pmatrix}
    =
    \begin{pmatrix}
      -1+1& -3-(-7)\\
      -4  &0
    \end{pmatrix}
    =
    \begin{pmatrix}
      0& 4\\
      -4  &0
    \end{pmatrix}\\    
    \transposed{A}&=
    \begin{pmatrix}
      1-a& 4\\
      3+b &0
    \end{pmatrix}
    =
      \begin{pmatrix}
      1-1& 4\\
      3+-7 &0
    \end{pmatrix}
    =
      \begin{pmatrix}
      0& 4\\
      -4 &0
    \end{pmatrix}
\end{align*}
となり$A$は交代行列である.  
\end{answerof}

\begin{answerof}{prob:1:6}
  \begin{align*}
    \begin{pmatrix}
      10\\4
    \end{pmatrix}
    +
    \begin{pmatrix}
      -2\\-3
    \end{pmatrix}
    +
    \begin{pmatrix}
      -3\\6
    \end{pmatrix}
    =
    \begin{pmatrix}
      10-2-3\\4-3+6
    \end{pmatrix}
    =
    \begin{pmatrix}
      5\\7
    \end{pmatrix}.
  \end{align*}

  \begin{align*}
    \begin{pmatrix}
      6&0\\0&4
    \end{pmatrix}+
    \begin{pmatrix}
      0&3\\-7&0
    \end{pmatrix}
    =
    \begin{pmatrix}
      6+0&0+3\\0-7&4+0
    \end{pmatrix}
    =
    \begin{pmatrix}
      6&3\\-7&4
    \end{pmatrix}.
  \end{align*}

  \begin{align*}
    \begin{pmatrix}
      -9&2\\3&7
    \end{pmatrix}
    \begin{pmatrix}
      1&0\\-5&4
    \end{pmatrix}=
    \begin{pmatrix}
     -9\cdot 1+2\cdot (-5)&-9\cdot 0+2\cdot 4\\3\cdot 1+7\cdot (-5)&3\cdot 0 +7\cdot 4
    \end{pmatrix}=
    \begin{pmatrix}
     -19&8\\-32&28
    \end{pmatrix}.
  \end{align*}

  \begin{align*}
    \begin{pmatrix}
      -9&2\\3&7
    \end{pmatrix}
    \begin{pmatrix}
      1\\-5
    \end{pmatrix}=
    \begin{pmatrix}
     -9\cdot 1+2\cdot (-5)\\3\cdot 1+7\cdot (-5)
    \end{pmatrix}=
    \begin{pmatrix}
     -19\\-32
    \end{pmatrix}.
  \end{align*}


  
  \begin{align*}
    \begin{pmatrix}
      -9&2
    \end{pmatrix}
    \begin{pmatrix}
      1\\-5
    \end{pmatrix}
    =
    \begin{pmatrix}
      (-9)\cdot 1 + 2\cdot (-5)
    \end{pmatrix}
    =
    \begin{pmatrix}
      -19
    \end{pmatrix}.
  \end{align*}

  \begin{align*}
    \begin{pmatrix}
      1\\-5
    \end{pmatrix}
    \begin{pmatrix}
      -9&2
    \end{pmatrix}
    =
    \begin{pmatrix}
      1\cdot (-9)& 1\cdot 2\\-5\cdot 1&-5\cdot 2
    \end{pmatrix}
    =
    \begin{pmatrix}
      -9& 2\\-5&-10
    \end{pmatrix}.
  \end{align*}

  \begin{align*}
    \begin{pmatrix}
      1&-3\\2&5
    \end{pmatrix}^2
    =
    \begin{pmatrix}
      1&-3\\2&5
    \end{pmatrix}
    \begin{pmatrix}
      1&-3\\2&5
    \end{pmatrix}
    =
    \begin{pmatrix}
      1\cdot 1+(-3)\cdot 2&1\cdot(-3)+(-3)\cdot 5\\2\cdot1 +5\cdot2&2\cdot(-3) +5\cdot5
    \end{pmatrix}
    =
    \begin{pmatrix}
      -5&-18\\12&19
    \end{pmatrix}.
  \end{align*}
  この結果は次の様に検算をすることができる:
    \begin{align*}
    A=
    \begin{pmatrix}
      1&-3\\2&5
    \end{pmatrix}
  \end{align*}
    とおくと,
  \cref{thm:cht:2dim}から,
  \begin{align*}
    A^2-(1+5)A+(1\cdot 5-(-3)\cdot2)E_2&=O_{2,2}\\
    A^2-6A+11E_2&=O_{2,2}\\
    A^2&=6A-11E_2.
  \end{align*}
  となる.
  したがって,
  \begin{align*}
    A^2&=6A-11E_2=
   6\begin{pmatrix}
      1&-3\\2&5
   \end{pmatrix}
   -11
  \begin{pmatrix}
      1&0\\0&1
  \end{pmatrix}
  =
  \begin{pmatrix}
      6-11&-18\\12&30-11
   \end{pmatrix}
  =
  \begin{pmatrix}
      -5&-18\\12&19
   \end{pmatrix}.
  \end{align*}
\end{answerof}

\begin{answerof}{prob:1:7}
  \begin{align*}
    A=\begin{pmatrix}
      9&0\\0&5
    \end{pmatrix}
  \end{align*}
  とし,
  \begin{align*}
    \begin{pmatrix}
      a_n&b_n\\c_n&d_n
    \end{pmatrix}
    =A^n
  \end{align*}
  とする.
  このとき,
  \begin{align*}
    \begin{pmatrix}
      a_{n+1}&b_{n+1}\\c_{n+1}&d_{n+1}
    \end{pmatrix}
    =A^{n+1}
    =AA^{n}
    =
    \begin{pmatrix}
      9&0\\0&5
    \end{pmatrix}
    \begin{pmatrix}
      a_n&b_n\\c_n&d_n
    \end{pmatrix}
    =
    \begin{pmatrix}
      9a_n+0c_n&9b_n+0d_n\\0a_n+5c_n&0b_n+5d_n
    \end{pmatrix}
    =
    \begin{pmatrix}
      9a_n&9b_n\\5c_n&5d_n
    \end{pmatrix}
  \end{align*}
  となるので,
  \begin{align*}
    \begin{pmatrix}
      a_{n+1}&b_{n+1}\\c_{n+1}&d_{n+1}
    \end{pmatrix}
    =
    \begin{pmatrix}
      9a_n&9b_n\\5c_n&5d_n
    \end{pmatrix}.
  \end{align*}
  よって,
  \begin{align*}
    a_n&=
    \begin{cases}
      9a_{n-1} &(n>1)\\
      9&(n=1)
    \end{cases}\\
    b_n&=
    \begin{cases}
      9b_{n-1}&(n>1)\\
      0&(n=1)
    \end{cases}\\
    c_n&=
    \begin{cases}
      5c_{n-1}&(n>1)\\
      0&(n=1)
    \end{cases}\\
    d_n&=
    \begin{cases}
      5d_{n-1}&(n>1)\\
      5&(n=1)
    \end{cases}
  \end{align*}
  となるので,
  \begin{align*}
    a_n&=9^{n-1}\cdot 9=9^n,\\
    b_n&=9^{n-1}\cdot 0=0,\\
    c_n&=5^{n-1}\cdot 0=0,\\
    d_n&=5^{n-1}\cdot 5=5^n.
  \end{align*}
  したがって,
  \begin{align*}
    A^n=\begin{pmatrix}9^n&0\\0&5^n\end{pmatrix}.
  \end{align*}

  この結果は次のように検算できる:
  $n=1$のとき
\begin{align*}
  A&=\begin{pmatrix}9&0\\0&5\end{pmatrix}\\
  \begin{pmatrix}9^n&0\\0&5^n\end{pmatrix}&=\begin{pmatrix}9&0\\0&5\end{pmatrix}.
\end{align*}
となり正しい.
  $n=2$のとき
\begin{align*}
  A^2&=\begin{pmatrix}9&0\\0&5\end{pmatrix}\begin{pmatrix}9&0\\0&5\end{pmatrix}
    =\begin{pmatrix}9\cdot 9+0\cdot 0&9\cdot 0+0\cdot5 \\0\cdot 9+5\cdot 0&0\cdot 0+5\cdot 5\end{pmatrix}
    =\begin{pmatrix}81&0\\0&25\end{pmatrix}
  \\
  \begin{pmatrix}9^2&0\\0&5^2\end{pmatrix}&=\begin{pmatrix}81&0\\0&25\end{pmatrix}.
\end{align*}
となり正しい.  
\end{answerof}


\begin{answerof}{prob:1:8}
  \begin{align*}
    A=
    \begin{pmatrix}
      1&2\\0&3
    \end{pmatrix}
  \end{align*}
  とし,
  \begin{align*}
    \begin{pmatrix}
      a_n&b_n\\c_n&d_n
    \end{pmatrix}
    =A^n
  \end{align*}
  とする.
  このとき,
  \begin{align*}
    \begin{pmatrix}
      a_{n+1}&b_{n+1}\\c_{n+1}&d_{n+1}
    \end{pmatrix}
    =A^{n+1}
    =AA^{n}
    =
    \begin{pmatrix}
      1&2\\0&3
    \end{pmatrix}
    \begin{pmatrix}
      a_n&b_n\\c_n&d_n
    \end{pmatrix}
    =
    \begin{pmatrix}
      1a_n+2c_n&1b_n+2d_n\\0a_n+3c_n&0b_n+3d_n
    \end{pmatrix}
    =
    \begin{pmatrix}
      a_n+2c_n&b_n+2d_n\\3c_n&3d_n
    \end{pmatrix}
  \end{align*}
  となるので,
  \begin{align*}
    \begin{pmatrix}
      a_{n+1}&b_{n+1}\\c_{n+1}&d_{n+1}
    \end{pmatrix}
    =
    \begin{pmatrix}
      a_n+2c_n&b_n+2d_n\\3c_n&3d_n
    \end{pmatrix}.
  \end{align*}
  よって,
  \begin{align*}
    a_n&=
    \begin{cases}
      a_{n-1}+2c_{n-1} &(n>1)\\
      1&(n=1)
    \end{cases}\\
    b_n&=
    \begin{cases}
      b_{n-1}+2d_{n-1}&(n>1)\\
      2&(n=1)
    \end{cases}\\
    c_n&=
    \begin{cases}
      3c_{n-1}&(n>1)\\
      0&(n=1)
    \end{cases}\\
    d_n&=
    \begin{cases}
      3d_{n-1}&(n>1)\\
      3&(n=1)
    \end{cases}
  \end{align*}
  となる. したがって,
  \begin{align*}
    c_n&=3^{n-1}\cdot 0\\
    d_n&=3^{n-1}\cdot 3=3^n
  \end{align*}
  である.
  したがって,
  \begin{align*}
    a_n&=
    \begin{cases}
      a_{n-1}+2c_{n-1} &(n>1)\\
      1&(n=1)
    \end{cases}\\
    &=
    \begin{cases}
      a_{n-1} &(n>1)\\
      1&(n=1)
    \end{cases}\\
    b_n&=
    \begin{cases}
      b_{n-1}+2d_{n-1}&(n>1)\\
      2&(n=1)
    \end{cases}\\
  &=
    \begin{cases}
      b_{n-1}+2\cdot 3^{n-1}&(n>1)\\
      2&(n=1)
    \end{cases}
  \end{align*}
  となる.
  したがって,
  \begin{align*}
    a_n&=1^{n-1}\cdot 1=1
  \end{align*}
  となる. また,
  \begin{align*}
    b_n
    &=2\cdot 3^{n-1}+b_{n-1}\\
    &=2\cdot 3^{n-1}+2\cdot 3^{n-2}+b_{n-2}\\
    &=\cdots\\
    &=2\cdot 3^{n-1}+2\cdot 3^{n-2}+b_{n-2}+\cdots +2\cdot 3^{2}+b_2\\
    &=2\cdot 3^{n-1}+2\cdot 3^{n-2}+b_{n-2}+\cdots +2\cdot 3^{2}+2\cdot 3^1 +b_1\\
    &=2\cdot 3^{n-1}+2\cdot 3^{n-2}+b_{n-2}+\cdots +2\cdot 3^{2}+2\cdot 3^1 +2\\
    &=2(3^{n-1}+3^{n-2}+\cdots+3+1)\\
    &=2\frac{3^{n}-1}{3-1}\\
    &=2\frac{3^{n}-1}{2}\\
    &=3^{n}-1.
  \end{align*}
  したがって,
  \begin{align*}
    A^n=\begin{pmatrix}1&3^{n}-1\\0&3^n\end{pmatrix}.
  \end{align*}


  この結果は次のように検算できる:
  $n=1$のとき,
  \begin{align*}
    A^n=
    \begin{pmatrix}
      1&2\\0&3
    \end{pmatrix}\\  
    \begin{pmatrix}1&2\cdot 3^{n}-1\cdot 2^{n+1}\\0&3^n\end{pmatrix}
      =
      \begin{pmatrix}1&2\cdot 3^{1}-1\cdot 2^{2}\\0&3^1\end{pmatrix}
        =
    \begin{pmatrix}1&2\\0&3\end{pmatrix}      .
  \end{align*}
  となり正しい.
  \begin{align*}
    AA^n&=
    \begin{pmatrix}
      1&2\\0&3
    \end{pmatrix}
    \begin{pmatrix}1&3^{n}-1\\0&3^n\end{pmatrix}
      =
    \begin{pmatrix}1\cdot 1+2\cdot 0&1\cdot(3^{n}-1)+2\cdot 3^n\\0\cdot 1+3\cdot 0&0\cdot (3^{n}-1)+3\cdot 3^n\end{pmatrix}
      =
    \begin{pmatrix}1\cdot 1+2\cdot 0&1\cdot(3^{n}-1)+2\cdot 3^n\\0\cdot 1+3\cdot 0&0\cdot (3^{n}-1)+3\cdot 3^n\end{pmatrix}
      =
    \begin{pmatrix}1&3^{n}-1+ 2\cdot 3^n\\3^{n+1}\end{pmatrix}\\
      =
    \begin{pmatrix}1&3\cdot 3^{n}-1\\3^{n+1}\end{pmatrix}\\
      =
    \begin{pmatrix}1&3^{n+1}-1\\3^{n+1}\end{pmatrix}\\
    AA^{n+1}&=\begin{pmatrix}1&3^{n+1}-1\\0&3^{n+1}\end{pmatrix}.
  \end{align*}
  となり等しい.
\end{answerof}

\begin{answerof}{prob:1:9}
 \Cref{thm:cht:2dim}から,
  \begin{align*}
    A^2-(1+5)A+(1\cdot 5-(-3)\cdot2)E_2&=O_{2,2}\\
    A^2-6A+11E_2&=O_{2,2}\\
    A^2&=6A-11E_2.
  \end{align*}
  したがって
  \begin{align*}
    A^3=A(6A-11E_2)=6A^2-11A=6(6A-11E_2)-11A=25A-66E_2.
  \end{align*}
  よって
  \begin{align*}
    A^3+A^2+A=(25A-66E_2)+(6A-11E_2)+A=32A-77E_2
    =32\begin{pmatrix}1&-3\\2&5\end{pmatrix}-77\begin{pmatrix}1&0\\0&1\end{pmatrix}\\
    =\begin{pmatrix}32-77&32\cdot(-3)\\32\cdot 2&32\cdot 5-77\end{pmatrix}
    =\begin{pmatrix}-45&-96\\64&83\end{pmatrix}
  \end{align*}

  また次のようにも計算できる:
  \begin{align*}
    \begin{pmatrix}
      1&-3\\2&5
    \end{pmatrix}^2
    =
    \begin{pmatrix}
      1&-3\\2&5
    \end{pmatrix}
    \begin{pmatrix}
      1&-3\\2&5
    \end{pmatrix}
    =
    \begin{pmatrix}
      1\cdot 1+(-3)\cdot 2&1\cdot(-3)+(-3)\cdot 5\\2\cdot1 +5\cdot2&2\cdot(-3) +5\cdot5
    \end{pmatrix}
    =
    \begin{pmatrix}
      -5&-18\\12&19
    \end{pmatrix}.
  \end{align*}

  \begin{align*}
    \begin{pmatrix}
      1&-3\\2&5
    \end{pmatrix}^3
    =
    \begin{pmatrix}
      1&-3\\2&5
    \end{pmatrix}
    \begin{pmatrix}
      1&-3\\2&5
    \end{pmatrix}^2
    =
    \begin{pmatrix}
      1&-3\\2&5
    \end{pmatrix}
    \begin{pmatrix}
      -5&-18\\12&19
    \end{pmatrix}
    =
    \begin{pmatrix}
      1\cdot(-5)+(-3)\cdot 12&1\cdot(-18)+(-3)\cdot 19\\2\cdot(-5) +5\cdot12&2\cdot(-18) +5\cdot19
    \end{pmatrix}
    =
    \begin{pmatrix}
      -41&-75\\50&59
    \end{pmatrix}.
  \end{align*}
  \begin{align*}
    A^3+A^2+A=
    \begin{pmatrix}
      -41&-75\\50&59
    \end{pmatrix}
    +
    \begin{pmatrix}
      -5&-18\\12&19
    \end{pmatrix}
    +
    \begin{pmatrix}
      1&-3\\2&5
    \end{pmatrix}
    =
    \begin{pmatrix}
      -41-5+1&-75-18-3\\50+12+2&59+19+5
    \end{pmatrix}
    =
    \begin{pmatrix}
      -45&-96\\64&83
    \end{pmatrix}.    
  \end{align*}
\end{answerof}


\begin{answerof}{prob:1:10}
  \begin{align*}
    a\begin{pmatrix}
      1\\1
    \end{pmatrix}
    +b\begin{pmatrix}
      -1\\0
    \end{pmatrix}
    =
    \begin{pmatrix}
      a-b\\a
    \end{pmatrix}
  \end{align*}
  であるので,
  \begin{align*}
    \begin{pmatrix}
      -3\\1
    \end{pmatrix}
    =
    a\begin{pmatrix}
      1\\1
    \end{pmatrix}
    +b\begin{pmatrix}
      -1\\0
    \end{pmatrix}
  \end{align*}
  とすると,
  \begin{align*}
    \begin{cases}
      a-b=-3\\a=1
    \end{cases}
  \end{align*}
  を得る.
  したがって
  $a=1$, $b=4$を得る.
  実際,
  $(a,b)=(1,4)$とすると,
  \begin{align*}
    a\begin{pmatrix}
      1\\1
    \end{pmatrix}
    +b\begin{pmatrix}
      -1\\0
    \end{pmatrix}
    =
    \begin{pmatrix}
      1\\1
    \end{pmatrix}
    +4\begin{pmatrix}
      -1\\0
    \end{pmatrix}
    =
    \begin{pmatrix}
      -3\\1
    \end{pmatrix}
  \end{align*}
  となる.
\end{answerof}

\section{}

\begin{answerof}{quiz:2:1}
  \begin{align*}
    \det(\begin{pmatrix}
      -3&-5\\1&2
    \end{pmatrix})=-3\cdot 2 -(-5)\cdot 1=-6+5=-1
  \end{align*}
\end{answerof}


\begin{answerof}{quiz:2:2}
  \begin{align*}
    \det(
    \det(\begin{pmatrix}
      -x&5\\1&2-x
    \end{pmatrix})=-x(2-x)-5=x^2-2x-5
  \end{align*}
  であるので
  \begin{align*}
    \det(
    \det(\begin{pmatrix}
      -x&5\\1&2-x
    \end{pmatrix})=0
  \end{align*}
  とすると,
  \begin{align*}
    x^2-2x-5&=0\\
    (x-1)^2-\sqrt{5}^2&=0\\
    (x-1+\sqrt5)(x-1-\sqrt5)=0
  \end{align*}
  である.
  よって, $x\in \Set{x-1+\sqrt5, x-1-\sqrt5}$.
\end{answerof}


\begin{answerof}{quiz:2:3}
  \begin{align*}
    \begin{pmatrix}
      1&0\\0&1
    \end{pmatrix}
    \begin{pmatrix}
      1&0\\0&1
    \end{pmatrix}=
    \begin{pmatrix}
      1&0\\0&1
    \end{pmatrix}
    =E_2
  \end{align*}
  であるので,
  \begin{align*}
      \begin{pmatrix}
      1&0\\0&1
    \end{pmatrix}
  \end{align*}
  の逆行列は
  \begin{align*}
      \begin{pmatrix}
      1&0\\0&1
    \end{pmatrix}
  \end{align*}
  である.
\end{answerof}

\begin{answerof}{quiz:2:4}
  \begin{align*}
    A&=\begin{pmatrix}
      1&3\\4&2
    \end{pmatrix},&
    \bbb&=\begin{pmatrix}
      8\\1
    \end{pmatrix}.   
  \end{align*}
  とする.
\begin{enumerate}
\item
  \begin{align*}
    A^{-1}&=\frac{1}{1\cdot 2 - 3\cdot 4}\begin{pmatrix}
      2&-3\\-4&1
    \end{pmatrix}
    =
    \frac{1}{-10}\begin{pmatrix}
      2&-3\\-4&1
    \end{pmatrix}.
  \end{align*}
  
  これは次のように検算できる.
  求めた行列ともとの行列をかけると
  \begin{align*}
    \frac{1}{-10}\begin{pmatrix}
      2&-3\\-4&1
    \end{pmatrix}
    \begin{pmatrix}
      1&3\\4&2
    \end{pmatrix}
    =
    \frac{1}{-10}\begin{pmatrix}
      2\cdot 1-3\cdot 4&2\cdot 3-3\cdot 2\\
      -4\cdot 1+1\cdot 4&-4\cdot 3+1\cdot2
    \end{pmatrix}
    =
    \frac{1}{-10}\begin{pmatrix}
      -10&0\\
      0&-10
    \end{pmatrix}
    =
    \begin{pmatrix}
      1&0\\
      0&1
    \end{pmatrix}
  \end{align*}
  となり, $E_2$となる.

\item
  $A$が正則なので解は$A^{-1}\bbb$である.
  したがって,
  \begin{align*}
    \frac{1}{-10}\begin{pmatrix}
      2&-3\\-4&1
    \end{pmatrix}
    \begin{pmatrix}
      8\\1
    \end{pmatrix}
    =
    \frac{1}{-10}\begin{pmatrix}
      2\cdot 8 -3\cdot 1\\
      -4\cdot 8 +1\cdot 1\\
    \end{pmatrix}
    =
    \frac{1}{-10}\begin{pmatrix}
      13\\-31
    \end{pmatrix}
    =
    \begin{pmatrix}
      \frac{-13}{10}\\\frac{31}{10}
    \end{pmatrix}
  \end{align*}

  これは次のように検算できる.
  \begin{align*}
    A\xx=
    \begin{pmatrix}
      1&3\\4&2
    \end{pmatrix}
    \begin{pmatrix}
      \frac{-13}{10}\\\frac{31}{10}
    \end{pmatrix}
    =
    \begin{pmatrix}
      \frac{-13}{10}+3\frac{31}{10}\\4\frac{-13}{10}+2\frac{31}{10}
    \end{pmatrix}
    =
    \begin{pmatrix}
      \frac{-13}{10}+\frac{93}{10}\\\frac{-52}{10}+2\frac{62}{10}
    \end{pmatrix}
    =
    \begin{pmatrix}
      8\\1
    \end{pmatrix}
  \end{align*}
  となり$\bbb$となる.
\end{enumerate}
\end{answerof}
